% examples/tikz-standalone.tex
%
% Smoke test for tikzlibrarykyoto.code.tex used completely on its own --
% no beamer, no `\usetheme{Kyoto}', no beamer...themeKyoto.sty file of
% any kind. This is what README.md means when it says the TikZ library
% "works standalone": with no Kyoto Beamer theme loaded, the library
% falls back to its own copies of the Kyoto colors (see
% tikzlibrarykyoto.code.tex's `\@ifundefinedcolor' guards) instead of
% reusing colors the theme would otherwise have defined.
%
% Build from the repository root with:
%
%   make
%
% which runs latexmk (LuaLaTeX) via examples/latexmkrc for all three
% example PDFs -- no need to install anything into TEXMF or set
% TEXINPUTS by hand. Equivalently, from this directory:
% `latexmk tikz-standalone.tex'.
%
\documentclass[tikz]{standalone}
\usetikzlibrary{kyoto}

\begin{document}
\begin{tikzpicture}[node distance=2cm]
  \node[kyoto/box] (a) {A};
  \node[kyoto/box outline alert, right=of a] (b) {B};
  \draw[kyoto/arrow] (a) -- (b);

  \node[
    kyoto/callout outline,
    callout absolute pointer={(b.north)},
    above=1cm of b,
  ] {Standalone};

  \begin{scope}[on background layer]
    \node[kyoto/highlight, fit=(a)(b)] {};
  \end{scope}
\end{tikzpicture}
\end{document}
